
which we shall denote by
@@M0@@

\sgaStableTarget{sga6.e10.ex.4-2}\textbf{Example 4.2.} Suppose that the scheme @@M1@@ has dimension @@M2@@. Take an element @@M3@@, for example @@M4@@, where @@M5@@ is an integral subscheme of dimension @@M6@@ in @@M7@@. Also take a homogeneous polynomial of weight @@M8@@
@@M9@@
where @@M10@@ is assigned weight @@M11@@ @@M12@@, for example a monomial @@M13@@ with @@M14@@. Substituting in @@M15@@, for each variable @@M16@@ @@M17@@, the Chern class
@@M18@@
gives an element
@@M19@@
The integer
@@M22@@
is called the \emph{Chern number of} @@M20@@ \emph{relative to} @@M21@@.
When @@M23@@, the reference to @@M24@@ is omitted from this terminology.

\sgaStableTarget{sga6.e10.sec.4.3}\subsection*{4.3.}
Without any hypothesis on the dimension of @@M25@@, consider a sequence @@M26@@ of invertible @@M27@@-modules and a coherent @@M28@@-module @@M29@@ such that @@M30@@. Put @@M31@@ @@M32@@ and @@M33@@.

